\documentclass[11pt,a4paper]{article}

% ===== XeLaTeX + Korean =====
\usepackage{kotex}
\usepackage{fontspec}
\setmainfont{Times New Roman}
\setmainhangulfont{AppleMyungjo}
\setsanshangulfont{Apple SD Gothic Neo}

% ===== Layout =====
\usepackage[a4paper,top=18mm,bottom=18mm,left=20mm,right=20mm]{geometry}
\usepackage{fancyhdr}
\setlength{\headheight}{24pt}
\usepackage{enumitem}
\usepackage{mdframed}
\usepackage{array}

\setlength{\parindent}{1em}
\linespread{1.18}

% ===== Header / Footer =====
\pagestyle{fancy}
\fancyhf{}
\renewcommand{\headrulewidth}{0.6pt}
\fancyhead[L]{\sffamily\bfseries 2026 고1 영어 변형 — 지문 26}
\fancyhead[R]{\sffamily George Bird Grinnell}
\renewcommand{\footrulewidth}{0pt}
\fancyfoot[C]{\framebox[16mm][c]{\thepage}}

% ===== helpers =====
\newcommand{\qno}[1]{\par\vspace{8pt}\noindent{\sffamily\bfseries #1.}\ }
\newcommand{\instr}[1]{{\sffamily #1}\par\vspace{3pt}}
\newlist{choices}{enumerate}{1}
\setlist[choices]{label=,leftmargin=1.5em,itemsep=2pt,topsep=3pt,parsep=0pt}

\newmdenv[linewidth=0.6pt,roundcorner=3pt,innertopmargin=7pt,innerbottommargin=7pt,
          innerleftmargin=9pt,innerrightmargin=9pt,skipabove=5pt,skipbelow=5pt]{pbox}
\newcommand{\nemo}[1]{\fbox{\,#1\,}}

\begin{document}

\begin{center}
{\fontsize{15}{17}\selectfont\sffamily\bfseries [지문 26] 다음 글을 읽고 물음에 답하시오.}
\end{center}
\vspace{4pt}

% ===== 공통 지문 (원본 단어 유지 + 어법 네모) =====
\begin{pbox}
\hspace{1em}George Bird Grinnell, one of the leading figures in early American conservation, was born in New York in 1849. Growing up near Audubon Park, he (A)\,\nemo{developed / was developing} an early interest in birds. After graduating from Yale University, he joined his first scientific expedition,* (B)\,\nemo{which / that} led to his career in American conservation. From the late 1870s, he wrote many articles in \textit{Forest and Stream} calling for wildlife conservation. Grinnell contributed to founding early conservation organizations, including the Audubon Society, and helped with the establishment of Glacier National Park in 1910. Through his writings and public campaigns, he had an influence on making early wildlife protection laws. Even after his death in 1938, his work in conservation (C)\,\nemo{continues / is continued} to inspire conservationists.

\vspace{4pt}
\noindent\footnotesize{* expedition: 탐사}
\end{pbox}

% ===================== Q1 =====================
\qno{1}\instr{(A), (B), (C)의 각 네모 안에서 어법에 맞는 표현으로 가장 적절한 것은?}
\begin{center}
\renewcommand{\arraystretch}{1.25}
\begin{tabular}{c c c c}
 & (A) & (B) & (C) \\
① & developed & which & continues \\
② & developed & which & is continued \\
③ & developed & that & continues \\
④ & was developing & which & continues \\
⑤ & was developing & that & is continued \\
\end{tabular}
\end{center}

% ===================== Q2 =====================
\qno{2}\instr{윗글의 내용과 일치하지 \underline{않는} 것은?}
\begin{choices}
\item ① Grinnell spent his formative years in the vicinity of a park named after a famous naturalist.
\item ② He participated in a scientific expedition prior to completing his studies at Yale University.
\item ③ His written contributions to a natural history journal advocated for the protection of wild animals.
\item ④ He was involved in establishing a national park that was officially created in the early twentieth century.
\item ⑤ His influence on conservation policy did not end with his passing in the late 1930s.
\end{choices}

% ===================== Q3 (서술형) =====================
\qno{3}\instr{[서술형] 윗글에 따르면, Grinnell이 야생동물 보호에 기여한 방식을 두 가지 이상 우리말로 쓰시오. (각 항목은 간결하게, 완전한 문장 불필요)}

\vspace{4pt}
\noindent $\Rightarrow$ \rule{10cm}{0.4pt}

\vspace{8pt}
\noindent \rule{\linewidth}{0.4pt}

% ===================== Q4 =====================
\qno{4}\instr{윗글의 제목으로 가장 적절한 것은?}
\begin{choices}
\item ① George Bird Grinnell: A Voice for Early Wildlife Conservation
\item ② The Founding of Yale University's First Scientific Expedition
\item ③ Why Glacier National Park Became America's Largest Park
\item ④ The Decline of Conservation Writing in the Late Nineteenth Century
\item ⑤ Audubon Park and the Rise of Modern City Planning
\end{choices}

\end{document}
